\section{Discussion}
\label{sec:discussion}

\paragraph{Scope of transfer results.}
The transfer experiment (Section~\ref{sec:transfer}) covers only the
\textbf{digit subset} (10 of 46 DHCD classes).  Extending transfer to the
full 36-consonant inventory is not currently possible: the only publicly
available Devanagari handwriting collections that share DHCD's consonant
label taxonomy — principally the Acharya~\cite{acharya2015dhcd} and
Pant datasets — are same-provenance splits of DHCD itself, and using them
as held-out targets would reintroduce contamination between training and
evaluation.  Public literature does not yet provide genuinely independent
consonant-level benchmarks, so transferability claims are limited to
digits.  This constraint itself underscores a broader infrastructure gap:
the Devanagari recognition community currently lacks the diversity of
held-out corpora that would allow systematic transfer studies, and the
digit result — where the compact model outperforms a supervised baseline
by \textbf{+14\,pp} under domain shift — suggests such benchmarks would
be worth building.

\paragraph{Single-benchmark evaluation.}
All consonant-level accuracy results in this paper are reported on a
single held-out split of DHCD.  The McNemar test controls for
within-split variability via paired comparison, but without a second,
independent consonant benchmark the findings may not generalise beyond
this split.  The saturation conclusion depends on properties specific to
DHCD — its $32{\times}32$ resolution, its concentrated label noise, and
its writer pool — that another dataset might not share.  This is not a
weakness peculiar to this study; it reflects the current state of public
Devanagari corpora.  Readers should interpret the saturation claim as a
property of DHCD at this resolution and provenance, not necessarily of
Devanagari recognition as a problem class.

\paragraph{Efficiency versus accuracy.}
The \textbf{15.6$\times$} parameter reduction and \textbf{9.5$\times$}
latency improvement come at no measurable accuracy cost on DHCD at its
current saturation point.  This finding supports a general principle: on
saturated benchmarks, compactness is not a compromise but a rational
choice, since the errors that remain are inaccessible to any model
regardless of capacity.  Whether this efficiency advantage persists under
a more discriminative benchmark — one where larger models could exploit
additional capacity to push below the 11-error floor — is an open
question this work cannot answer.  These results therefore support
deployment on DHCD more directly than they support a general claim about
efficiency under harder conditions, and they motivate the construction of
such harder benchmarks as the most productive next step for the field.
